%%%%%%%%%%%%%%%%%%%%%%%%%%%%%%%%%%%%%%%%%%%%%%%%%%%%%%%%%%%%%%%%%%%%
% bare_jrnl.tex
% V1.4b
% 2015/08/26
% by Michael Shell
% see http://www.michaelshell.org/
% for current contact information.                                                                 
%%%%%%%%%%%%%%%%%%%%%%%%%%%%%%%%%%%%%%%%%%%%%%%%%%%%%%%%%%%%%%%%%%%%

%% This is a skeleton file demonstrating the use of IEEEtran.cls
%% (requires IEEEtran.cls version 1.8b or later) with an IEEE
%% journal paper.
%%
%% Support sites:
%% http://www.michaelshell.org/tex/ieeetran/
%% http://www.ctan.org/pkg/ieeetran
%% and
%% http://www.ieee.org/


\documentclass[journal, onecolumn]{IEEEtran}



%
\usepackage{instructivo}  % Commands for multiple choice
\graphicspath{{./}{./fig/}}

\usepackage[natbib=true,style=trad-unsrt,backend=biber]{biblatex}
\addbibresource{literatura.bib}

\usepackage{standalone}
\usepackage{float}
\usepackage{graphicx}
\usepackage{booktabs}
\usepackage{tabularx}
\usepackage{url}
\usepackage{hyperref}



\renewcommand\IEEEkeywordsname{Palabras clave}
% correct bad hyphenation here
\hyphenation{op-tical net-works semi-conduc-tor}

\begin{document}

%
% paper title
% Titles are generally capitalized except for words such as a, an, and, as,
% at, but, by, for, in, nor, of, on, or, the, to and up, which are usually
% not capitalized unless they are the first or last word of the title.
% Linebreaks \\ can be used within to get better formatting as desired.
% Do not put math or special symbols in the title.
\title{Primer proyecto de aplicación: amplificador de audio}


\author{Matías~A.~Camacho,~\IEEEmembership{Estudiante,~TEC}
        ,~Samantha G.~Cubero,~\IEEEmembership{Estudiante,~TEC}
        ~y~Andrés~Rojas,~\IEEEmembership{Estudiante,~TEC}
}

% The paper headers
\markboth{Taller de diseño analógico, IS2026}%
{Taller de diseño analógico}


% make the title area
\maketitle




\clearpage
\tableofcontents
\clearpage

% Se sugiere no más de cuatro palabras o frases cortas en orden alfabético, separadas por comas, que representen su reporte
\begin{IEEEkeywords}
Diodo, tensión, corriente, junta, ruptura.
\end{IEEEkeywords}


%%%%%%%%%%%%%%%%%%%%%%%%%%%%%%%%%%%%%%%%%%%%%%%%%%%%%%%%%%%%%%%%%%%%%%%%%%%%%%%%%%%%%%%%%%%%%%
%%%%%%%%%%%%%%%%%%%%%%%%%%%%%%%%%%%%%%%%%%%%%%%%%%%%%%%%%%%%%%%%%%%%%%%%%%%%%%%%%%%%%%%%%%%%%%
\section{Introducción}


\subsection{}
E
\subsection{}



\subsection{Especificaciones del diseño}


%%%%%%%%%%%%%%%%%%%%%%%%%%%%%%%%%%%%%%%%%%%%%%%%%%%%%%%%%%%%%%%%%%%%%%%%%%%%%%%%%%%%%%%%%%%%%%
%%%%%%%%%%%%%%%%%%%%%%%%%%%%%%%%%%%%%%%%%%%%%%%%%%%%%%%%%%%%%%%%%%%%%%%%%%%%%%%%%%%%%%%%%%%%%%
\section{Diseño}
El proyecto propone diseñar un circuito PWM solo con la utilización de una fuente de DC y varios circuitos con topologías de amplificadores operacionales. Específicamente, se planteó generar una señal sinusoidal con un puente Wien, para compararla (el circuito integrado LM393 funcionando como comparador) con una señal de triangular. La señal triangular se generaría con un circuito integrador alimentado con una señal cuadrada, la cual se generaría con un circuito comparador.
Como parámetros de diseño, se tiene que la onda sinusoidal debe tener una frecuencia de 100~Hz y magnitud de 2~V. La onda transportadora debe tener una magnitud de 2.1~V a 2.5~V y una frecuencia de 5~kHz.


\section{Simulaciones y análisis}




\section{Resultados Experimentales}





%%%%%%%%%%%%%%%%%%%%%%%%%%%%%%%%%%%%%%%%%%%%%%%%%%%%%%%%%%%%%%%%%%%%%%%%%%%%%%%%%%%%
%%%%%%%%%%%%%%%%%%%%%%%%%%%%%%%%%%%%%%%%%%%%%%%%%%%%%%%%%%%%%%%%%%%%%%%%%%%%%%%%%%%%
\section{Conclusiones}


%%%%%%%%%%%%%%%%%%%%%%%%%%%%%%%%%%%%%%%%%%%%%%%%%%%%%%%%%%%%%%%%%%%%%%%%%%%%%%%%%%%%
%%%%%%%%%%%%%%%%%%%%%%%%%%%%%%%%%%%%%%%%%%%%%%%%%%%%%%%%%%%%%%%%%%%%%%%%%%%%%%%%%%%%
%\appendices





%%%%%%%%%%%%%%%%%%%%%%%%%%%%%%%%%%%%%%%%%%%%%%%%%%%%%%%%%%%%%%%%%%%%%%%%%%%%%%%%%%%%
%%%%%%%%%%%%%%%%%%%%%%%%%%%%%%%%%%%%%%%%%%%%%%%%%%%%%%%%%%%%%%%%%%%%%%%%%%%%%%%%%%%%

\printbibliography

\vfill

\end{document}






